
\def\PUDcodigo{ACE016}

\if\COMPILAR0
   \def\PUDcodacad{04507.47}
\else
   \def\PUDcodacad{04506.46}
\fi

\def\PUDdisciplina{Projetos Sociais}

\def\PUDsemestre{S10}

\def\PUDhoras{40}

%NOVOS ---------------------------------------------------
\def\PUDhorasTeoria{20}
\def\PUDhorasPratica{20}

\def\PUDinterECA{---}

\def\PUDatuaisECA{---}

\def\PUDrecursos{Projetor multimídia; Quadro branco e pincel; Práticas de campo.}

%----------------------------------------------------------

\def\PUDcreditos{2}

\def\PUDprerequisito{---}

\def\PUDpreacad{---}

\def\PUDobjetivos{Compreender temáticas ligadas à cidadania no contexto contemporâneo brasileiro; conceituar projetos sociais;
estudar projetos sociais exemplares; conhecer e participar de ações e projetos sociais da comunidade local; elaborar e executar ações, projetos e programas sociais.}

\def\PUDementa{Cidadania, Sociedade Civil, Estado e Movimentos Sociais (minorias sociais, gênero, comunidades étnicas, tradicionais e populares, urbanas e rurais). Conceituação de Projetos Sociais. Estudos de casos exemplares. Elaboração de programas, projetos e ações sociais. Práticas em Projetos Sociais. }

\def\PUDconteudo{ &  HISTÓRIA DOS MOVIMENTOS SOCIAIS NO BRASIL CONTEMPORÂNEO: Cidadania – conceito e exercício social; os anos 1960/1970 e a perca dos direitos civis; os anos 1980 e a eclosão dos novos sujeitos sociais e suas práticas (negros, indígenas, imigrantes, mulheres, homossexuais, trabalhadores urbanos, trabalhadores rurais, bairros e favelas, comunidades tradicionais etc.); ONGs, Sociedade Civil e Estado no Brasil contemporâneo; ONGs e projetos Sociais. 
\\
 &  PROJETOS SOCIAIS: conceituação e terminologia afins; estudos de casos.
\\ 
 &  PRÁTICA EM PROJETOS SOCIAIS I: conhecimento de ONGs e Projetos Sociais da comunidade local; análise de ONGs e Projetos Sociais da comunidade local; planejamento e elaboração e Ações/Projetos Sociais para a comunidade local.
\\
 &  PRÁTICA EM PROJETOS SOCIAIS II: execução de Ações/Projetos Sociais na comunidade local; avaliação de Ações/Projetos Sociais na comunidade local. }

\def\PUDmetodo{Aulas expositivas; Seminários; Apresentação e discussão de artigos de jornais e/ou literatura especializada; Aulas de Campo; Visitas Técnicas; Práticas em Projetos Sociais. A Prática como Componente Curricular de Ensino poderá ser ministrada através de: aulas expositivas, criação e aplicação de técnicas de ensino, apresentação de seminários e elaboração de material didático. }

\def\PUDavalia{A avaliação será desenvolvida ao longo do semestre, de forma processual e contínua, valorizando os aspectos qualitativos em relação aos quantitativos, onde os critérios a serem avaliados serão: conhecimento individual sobre temas relativos aos assuntos estudados em sala; grau de participação do aluno em atividades que exijam produção individual e em
equipe; planejamento, organização, coerência de idéias e clareza na elaboração de trabalhos escritos ou destinados à demonstração do domínio dos conhecimentos técnico-pedagógicos e científicos adquiridos; criatividade e o uso de recursos diversificados; domínio de atuação discente (postura e desempenho). 

A avaliação da Prática como Componente Curricular seguirá os critérios citados anteriormente em conformidade com a metodologia estabelecida para a disciplina. Será avaliado também as ações/projetos elaborados e/ou executados pelos alunos. Ocorrerá também avaliação somativa de acordo com o Regulamento da Organização Didática (ROD) do IFCE. \vspace{12pt}}

\def\PUDbibbasica{ &  \href{www.biblioteca.ifce.edu.br/index.asp?codigo_sophia=61857}{Demo}, Pedro.  Participação é conquista.  6º ed.  São Paulo, SP: Cortez, 2009.  176p.  ISBN: 9788524901287.
%\\ 
 %&   AGUILAR, Maria José;  ANDER-EGG, Ezequiel.  Avaliação de Programas e Serviços Sociais.  Petrópolis: Vozes, 1994.  ISBN 9788532612199
%\\ 
 %&   CAPRILES, René.  Makarenko O Nascimento da Pedagogia Socialista.  São Paulo: Scipione, 1989.  ISBN 9788526244108. }
\\
 &  \href{www.biblioteca.ifce.edu.br/index.asp?codigo_sophia=62049}{Drucker}, Peter F.  Inovação e espírito empreendedor: prática e princípios.  São Paulo, SP: Cengage Learning, 2015.  378p.  ISBN: 9788522108596.
\\
 & \href{www.biblioteca.ifce.edu.br/index.asp?codigo_sophia=65379}{Salim}, C. S.  Introdução ao empreendedorismo: despertando a atitude empreendedora.  Rio de Janeiro: Elsevier, 2010.  245p.  ISBN: 9788535234664. }

\def\PUDbibcomplementar{ &  \href{www.biblioteca.ifce.edu.br/index.asp?codigo_sophia=61985}{Chiavenato}, Idalberto.  Empreendedorismo: dando asas ao espírito empreendedor.  4ª ed.  Barueri, SP: Manole, 2012.  315p.  ISBN: 9788520432778.
%&   DRUCKER, P.  E.  Administração de Organizações sem Fins Lucrativos: Princípios e Práticas.  São Paulo: Pioneira, 1995.  ISBN 9788522101900
\\ 
 &  \href{www.biblioteca.ifce.edu.br/index.asp?codigo_sophia=61755}{Richardson}, Roberto Jarry.  Pesquisa social: métodos e técnicas.  3º ed.  São Paulo, SP: Atlas, 2015.  334p.  ISBN: 9788522421114.
\\
 &  \href{www.biblioteca.ifce.edu.br/index.asp?codigo_sophia=61750}{Maximiano}, Antonio Cesar Amaru.  Administração para empreendedores.  2º ed.  São Paulo, SP: Pearson Prentice Hall, 2011.  240p.  ISBN: 9788576058762.
\\
 &  \href{www.biblioteca.ifce.edu.br/index.asp?codigo_sophia=56904}{Ramos}, Ieda Cristina Alves; de Moura, Paulo G. M.; Giehl, Pedro Roque; Gianezini, Miguelangelo; dos Santos, Andréa; de Borba, Carolina dos Anjos; da Silveira, Luciana Conceição Lemos.  Captação de recursos para projetos sociais.  E-book.  Intersaberes.  126p.  ISBN: 9788582124901.
\\
 &  \href{www.biblioteca.ifce.edu.br/index.asp?codigo_sophia=58587}{Giehl}, Pedro Roque; Webler, Darlene Arlete; Ramos, Ieda Cristina Alves; Silveira, Luciana Conceição Lemos da; Gianezini, Miguelangelo.  Elaboração de projetos sociais.  E-book.  1º ed.  Intersaberes.  180p.  ISBN: 9788544302729. 
\\
 & \href{www.biblioteca.ifce.edu.br/index.asp?codigo_sophia=67527}{Nunes}, Antônia Elisabeth da Silva Souza; Oliveira, Elias Vieira de.  Implementação das diretrizes curriculares para a educação das relações étnico - raciais e o ensino de História e Cultura Afro - Brasileira e Africana na educação profissional e tecnológica.  Brasília, DF: MEC / Secretaria de Educação Profissional e Tecnológica, 2008.  180p.
\\
 & \href{www.biblioteca.ifce.edu.br/index.asp?codigo_sophia=55695}{Lafer}, Celso.  A Internacionalização dos Direitos Humanos: constituição, racismo e relações internacionais.  E-book.  Manole.  148p.  ISBN: 9788520424292. }

\def\PUDautor{Juliana Brito}

\def\PUDdata{2017-08-11}

\def\PUDrevisao{3}

\def\PUDrevisor{Samuel Vieira Dias}

\def\PUDdatarevisao{2017-08-11}

\PUDcorpo
